\subsection{Where should the probe read?}
\label{sec:readlocation}
Adapting the probe to duplex speech requires choosing its layer,
position, and read time. The final layer feeds the output head, but
this does not ensure that its failure signal survives a change in
query distribution or input modality. We select L22 based on its
robustness in the text-input layer audit and representation-transfer
tests (\S\ref{sec:duplex}--\ref{sec:xmodal}); those results alone do
not establish the best readout for native audio serving.

The native feature vector concatenates the latest L22 tail state,
the mean of the tail states, and the mean input state, allowing the
probe to combine the current boundary state with pooled context.
The separate native layer-by-readout ablation
(Figure~\ref{fig:native-ablation}) supports aggregation at L22;
it does not show a mid-layer advantage for a single-token read.

We read these features when the generation call first returns a
listen-to-speak transition. This makes the gate available at the
local model's response boundary without waiting for a completed
answer. The call can already contain the first generated text unit,
so we call the decision point \emph{answer onset}. Further answer
chunks may expose additional failure information, at the cost of a
later handoff (Appendix~\ref{app:readtime}). We fit the native probe
and thresholds on features and answer labels from this serving regime
to match its audio input and generation behavior. This calibration
is separate from the representation-transfer study, which alone
does not provide a zero-shot native gate.
